\PassOptionsToPackage{unicode=true}{hyperref} % options for packages loaded elsewhere
\PassOptionsToPackage{hyphens}{url}
%
\documentclass[14pt,]{extarticle}
\usepackage{lmodern}
\usepackage{amssymb,amsmath}
\usepackage{ifxetex,ifluatex}
\usepackage{fixltx2e} % provides \textsubscript
\ifnum 0\ifxetex 1\fi\ifluatex 1\fi=0 % if pdftex
  \usepackage[T1]{fontenc}
  \usepackage[utf8]{inputenc}
  \usepackage{textcomp} % provides euro and other symbols
\else % if luatex or xelatex
  \usepackage{unicode-math}
  \defaultfontfeatures{Ligatures=TeX,Scale=MatchLowercase}
    \setmainfont[]{Times New Roman}
  \ifxetex
    \usepackage{xeCJK}
    \setCJKmainfont[]{Microsoft JhengHei}
  \fi
  \ifluatex
    \usepackage[]{luatexja-fontspec}
    \setmainjfont[]{Microsoft JhengHei}
  \fi
\fi
% use upquote if available, for straight quotes in verbatim environments
\IfFileExists{upquote.sty}{\usepackage{upquote}}{}
% use microtype if available
\IfFileExists{microtype.sty}{%
\usepackage[]{microtype}
\UseMicrotypeSet[protrusion]{basicmath} % disable protrusion for tt fonts
}{}
\IfFileExists{parskip.sty}{%
\usepackage{parskip}
}{% else
\setlength{\parindent}{0pt}
\setlength{\parskip}{6pt plus 2pt minus 1pt}
}
\usepackage{hyperref}
\hypersetup{
            pdfborder={0 0 0},
            breaklinks=true}
\urlstyle{same}  % don't use monospace font for urls
\usepackage[margin=1in]{geometry}
\usepackage{longtable,booktabs}
% Fix footnotes in tables (requires footnote package)
\IfFileExists{footnote.sty}{\usepackage{footnote}\makesavenoteenv{longtable}}{}
\setlength{\emergencystretch}{3em}  % prevent overfull lines
\providecommand{\tightlist}{%
  \setlength{\itemsep}{0pt}\setlength{\parskip}{0pt}}
\setcounter{secnumdepth}{0}
% Redefines (sub)paragraphs to behave more like sections
\ifx\paragraph\undefined\else
\let\oldparagraph\paragraph
\renewcommand{\paragraph}[1]{\oldparagraph{#1}\mbox{}}
\fi
\ifx\subparagraph\undefined\else
\let\oldsubparagraph\subparagraph
\renewcommand{\subparagraph}[1]{\oldsubparagraph{#1}\mbox{}}
\fi

% set default figure placement to htbp
\makeatletter
\def\fps@figure{htbp}
\makeatother


\date{}

\begin{document}

\hypertarget{week-06-ux4f5cux696dux56deux994bposner-spatial-cueing-task-builder-pavlovia}{%
\section{Week 06 作業回饋（Posner Spatial Cueing Task --- Builder +
Pavlovia）}\label{week-06-ux4f5cux696dux56deux994bposner-spatial-cueing-task-builder-pavlovia}}

\hypertarget{ux5442ux6770ux9a5b-94100}{%
\subsection{113825002 呂杰驛 ---
94／100}\label{ux5442ux6770ux9a5b-94100}}

\hypertarget{ux7e3dux8a55}{%
\subsubsection{總評}\label{ux7e3dux8a55}}

以\textbf{單一 Trial\_Routine} 搭配各 component 的 start/stop
時間參數來重現 Fixation → Cue → SOA → Target
的完整時序，屬於作業允許的精簡實作方式（assignment note: ``You may
combine Fixation + Cue + SOA\_Blank into fewer
routines''），設計上乾淨有效率。Code Component 同時提供 Python 與
JavaScript 版本、成功同步至 Pavlovia 並完成 pilot 執行，且附有詳盡的
\texttt{Task\ Intro\_readme.md.txt} 說明實驗設定與輸出欄位。小幅扣分在於
Flow 中未納入 ITI／Rest 路由，以及 \texttt{Stimulus\ File} 未包含
\texttt{soa} 欄位。

\hypertarget{ux5404ux9805ux8a55ux8a9e}{%
\subsubsection{各項評語}\label{ux5404ux9805ux8a55ux8a9e}}

\begin{longtable}[]{@{}lll@{}}
\toprule
\begin{minipage}[b]{0.09\columnwidth}\raggedright
指標\strut
\end{minipage} & \begin{minipage}[b]{0.09\columnwidth}\raggedright
得分\strut
\end{minipage} & \begin{minipage}[b]{0.09\columnwidth}\raggedright
評語\strut
\end{minipage}\tabularnewline
\midrule
\endhead
\begin{minipage}[t]{0.09\columnwidth}\raggedright
Experiment Structure\strut
\end{minipage} & \begin{minipage}[t]{0.09\columnwidth}\raggedright
16/20\strut
\end{minipage} & \begin{minipage}[t]{0.09\columnwidth}\raggedright
Intro → Trial\_Routine → Thanks 三個 routine
已涵蓋主要流程，Trial\_Routine 內以 component 時序（fixation
0--1.0s、cue 0.5--0.6s、target 1.0s+）重現 SOA 設計，精簡合理。但 Flow
中未放入 ITI 與 Rest/BreakTime routine（\texttt{BreakTime\_Routine}
已定義但未納入 Flow），未完全對應作業指定的 8 個
routine。\textbf{(-4)}\strut
\end{minipage}\tabularnewline
\begin{minipage}[t]{0.09\columnwidth}\raggedright
Conditions File\strut
\end{minipage} & \begin{minipage}[t]{0.09\columnwidth}\raggedright
13/15\strut
\end{minipage} & \begin{minipage}[t]{0.09\columnwidth}\raggedright
\texttt{Stimulus\ File\_Posner\ Spatial\ Cueing\ Task.xlsx} 含
\texttt{cue\_side,\ target\_side,\ validity,\ correct\_key} 四欄，8
valid + 2 invalid 達到 80/20 比例 ✓。惟未包含作業所列之 \texttt{soa}
欄位（雖然 SOA 由 component
時間控制，為一致性仍建議加入）。\textbf{(-2)}\strut
\end{minipage}\tabularnewline
\begin{minipage}[t]{0.09\columnwidth}\raggedright
Loop \& Randomization\strut
\end{minipage} & \begin{minipage}[t]{0.09\columnwidth}\raggedright
10/10\strut
\end{minipage} & \begin{minipage}[t]{0.09\columnwidth}\raggedright
\texttt{Trial\_Loop} nReps=4、loopType=random、連結 Excel，產生 40 試驗
✓。\strut
\end{minipage}\tabularnewline
\begin{minipage}[t]{0.09\columnwidth}\raggedright
Cue \& Target Positioning\strut
\end{minipage} & \begin{minipage}[t]{0.09\columnwidth}\raggedright
15/15\strut
\end{minipage} & \begin{minipage}[t]{0.09\columnwidth}\raggedright
Trial\_Routine 內使用 CodeComponent 依
\texttt{cue\_side}／\texttt{target\_side} 動態設定座標，兩個
Polygon（cue box）與 target 的位置皆隨 trial 更新，符合 Approach
B。\strut
\end{minipage}\tabularnewline
\begin{minipage}[t]{0.09\columnwidth}\raggedright
Keyboard Response\strut
\end{minipage} & \begin{minipage}[t]{0.09\columnwidth}\raggedright
10/10\strut
\end{minipage} & \begin{minipage}[t]{0.09\columnwidth}\raggedright
Respkey 元件限制
\texttt{\textquotesingle{}left\textquotesingle{},\ \textquotesingle{}right\textquotesingle{}}、force
end of routine、store correct 設定 \texttt{\$correct\_key} ✓，CSV 中
\texttt{Respkey.corr} 與 \texttt{Respkey.rt} 正常紀錄。\strut
\end{minipage}\tabularnewline
\begin{minipage}[t]{0.09\columnwidth}\raggedright
Pavlovia Deployment\strut
\end{minipage} & \begin{minipage}[t]{0.09\columnwidth}\raggedright
15/15\strut
\end{minipage} & \begin{minipage}[t]{0.09\columnwidth}\raggedright
已同步 Pavlovia（包含
\texttt{.js}、\texttt{-legacy-browsers.js}、\texttt{index.html}），提供連結
\texttt{https://pavlovia.org/jelu/Posner\_Spatial\_Cueing\_Task}，Code
Component 有 JS 對應版本。\strut
\end{minipage}\tabularnewline
\begin{minipage}[t]{0.09\columnwidth}\raggedright
Pilot Data\strut
\end{minipage} & \begin{minipage}[t]{0.09\columnwidth}\raggedright
10/10\strut
\end{minipage} & \begin{minipage}[t]{0.09\columnwidth}\raggedright
\texttt{data/Test001\_Posner\ Spatial\ Cueing\ Task.csv} 含 40
筆試驗、\texttt{Respkey.corr}、\texttt{Respkey.rt}、condition labels
齊全。\strut
\end{minipage}\tabularnewline
\begin{minipage}[t]{0.09\columnwidth}\raggedright
Documentation\strut
\end{minipage} & \begin{minipage}[t]{0.09\columnwidth}\raggedright
5/5\strut
\end{minipage} & \begin{minipage}[t]{0.09\columnwidth}\raggedright
\texttt{Task\ Intro\_readme.md.txt} 清楚記載試驗數目、時序（ITI 500 ms /
cue 100 ms / target 1500 ms）、cue validity 80/20、以及 CSV
輸出欄位對應（Accuracy=Respkey.corr, RT=Respkey.rt），文件品質佳。\strut
\end{minipage}\tabularnewline
\bottomrule
\end{longtable}

\hypertarget{ux6539ux9032ux5efaux8b70}{%
\subsubsection{改進建議}\label{ux6539ux9032ux5efaux8b70}}

\begin{enumerate}
\def\labelenumi{\arabic{enumi}.}
\tightlist
\item
  \textbf{加入 ITI 與 Rest routine（+4 分）：} 將
  \texttt{BreakTime\_Routine} 拉入
  Flow，搭配區塊迴圈可讓受試者在長時間任務中休息；並在試驗間加入獨立 ITI
  routine 使結構更對應作業規格。
\item
  \textbf{補上 \texttt{soa} 欄位（+2 分）：} 即使 SOA 由 component
  時間控制，將 \texttt{soa=0.5} 寫入 Excel
  可提升條件檔的可追溯性與線上版的一致性。
\end{enumerate}

\end{document}
